\scanpage{98}

Soit $K$ une extension \deleted{\unsure{gal.}} \deleted{tot.} réelle finie de $Q$, $L$ ext. quadratique tot. imaginaire de $K$, $q=p^{\unsure{r}}$ une puissance de nb premier. On cherche les $\lambda\in K$ \uline{entiers} tels que $\lambda\bar\lambda=q^i$, i.e. $|\lambda|=q^{i/2}$, \inserted{pour $i$ convenable $\in N$.} Ce sont des unités exc. $p$, et ils engendrent \deleted{le groupe des} un sous-groupe du groupe \deleted{\unsure{U des} unités exc. $p$, \unsure{tel le $\lambda$ est}} formé des $\unsure{\lambda'|\lambda}$ unités exc. $p$ tels que $\lambda\bar\lambda=q^i$ ($i\in Z$). Notons que \deleted{ce sont des unités \uline{absolues} exc. $p$. Donc \unsure{finalement il suffit en $\lambda$ d'écrire} $p^{i/2}u$, avec $u$} si $q$ est un carré de $K$, on peut écrire les $\lambda$ sous la forme $\lambda=q^{i/2}u$, où $u$ est une unité exc. $p$ telle que $u\bar u=1$, ce qui implique que $u$ est une \uline{unité absolue} exc. $p$. Les $u$ forment un groupe $U$, et l'hom
\[
U\unsure{\text{—}|\longrightarrow} Z^P\qquad u\unsure{\sim|\longmapsto}(v_{\mathfrak p}(u))_{\mathfrak p\in P}
\]
(où $P$ désigne l'ens des $\mathfrak p$ de $L$ sur $p$) est de noyau fini. D'après la prop. plus haut, son image est \inserted{contenue et} d'indice fini dans $(Z^{\unsure{P}})'$, \deleted{$P'$ est l'ens des $\mathfrak p$} formé des diviseurs de trace \inserted{rel. à $K$} nulle. Donc si pour $\forall$ $p_0$ \deleted{\unsure{de $K$}} idéal premier de $K$ sur $p$ qui \uline{splitte en deux} dans $L$, on choisit un $\mathfrak p$ au dessus, on trouve un isom. mod. groupes finis $U\simeq Z^Q$, $Q$ est l'ens des $p_0$

\mspage{97}
